\section{Discussion}\label{sec:discussion}
\subsection{What the Results Establish}
The most consistent practical finding is the separation between attacker recognition and useful routing service. In the seven-attacker static main comparison, DIS-FDM and both AEF variants identify a similar fraction of attackers but preserve markedly different downward delivery. The route-repair ablation strengthens the mechanistic interpretation: identifying an origin is insufficient if the implementation continues to propagate damaging transit-parent information. This extends the evaluation focus beyond the DIS mitigation studied in \cite{prajisha2025} and complements DAO-focused mitigation \cite{limsd2024}. It does not imply that either original method was intended to solve the exact combined threat implemented here.

The method does not provide the best classification outcome in every cell. In the held-out extension, Li-MSD's substantially smaller false-positive burden yields higher macro-F1 even where AEF preserves more service. Secure-RPL also avoids the mobile false-isolation burden in the main comparison while remaining less effective at the DAO-related service damage. These trade-offs should guide interpretation of ``improvement'': the result depends on whether the operational priority is route preservation, reliable attribution, low false isolation, or alert speed. A deployment recommendation requires a specified cost for false isolation; the present experiments do not estimate that cost for an industrial installation.

\subsection{Learning, Evidence Dependence, and Negative Findings}
The learning hypothesis is not supported by the benign mobility comparison. Repeated Q-learning decisions are driven by an anomaly signal assembled from the same observations used by the fixed rule. The reward can therefore reward consistency with an erroneous internal signal rather than consistency with a true attack label. Mobility changes preferred parents legitimately, which makes this limitation especially relevant. The observed churn suggests that per-operation alert cost and unique-node FPR should both be reported. It does not, without further factorial interventions, establish that reward design alone caused every false isolation.

These findings do not contradict Q-learning as a general method \cite{watkins1992}, or establish a ranking over learning-based RPL security such as QSec-RPL and iTRPL \cite{qsecrpl2023,itrpl2024}. Those methods use different attacks, states, rewards, and decision placement. Likewise, combining Dempster--Shafer reasoning and reinforcement learning is already represented in other network settings \cite{gyawali2021,jasti2024}. The defensible contribution here is the controlled examination of the specific RPL observations and handling policy, including where the learner fails to add value.

\subsection{Research and Applied Significance}
The research significance is an empirical account of how a local consistency observation, evidence persistence, and mitigation action affect different outcomes under a same-node DIS/DAO attack. The strongest descriptive service effect is associated with selective DAO handling, whereas the extra decision machinery does not consistently improve attribution. This distinction gives a more precise direction for subsequent algorithm design than a blanket claim that adaptation improves intrusion detection.

The applied significance is the diagnostic protocol: a high conditional PDR can coexist with a large population of silent clients. The scale failure demonstrates why a benign connectivity gate and a network-wide denominator must precede a large deployment comparison. Reporting actual scenario coordinates and stack warnings also makes the failure inspectable. These are findings about the evaluated implementation and simulator; they do not establish field reliability, production readiness, or successful 1000-node deployment.

\section{Limitations and Threats to Validity}\label{sec:limitations}
\textit{Internal validity.} A shared seed controls scenario generation, but defense decisions can change subsequent packet timing and random-number consumption inside the simulation. Paired runs are comparable realizations, not identical packet-level counterfactual histories. Configuration, timer, and baseline implementation choices can affect outcomes. The source audit excludes obvious privileged-label references but is not a formal proof against every information leak. Online calibration occurs during a deliberately benign interval, and mobility-specific thresholds are selected using known scenario regimes.

\textit{Construct validity.} Network-level any-observer detection is weaker than consistent correct attribution by all observers and may credit historical pre-onset flags. The extension's new-alert analysis mitigates, but does not eliminate, this interpretive difference. One persistent event per attacker does not represent repeated independent intrusions. A high downward PDR is conditional on reaching the root; it is therefore considered alongside participation, expected delivery, and ARI. Radio power, native object sizes, and callback clocks do not substitute for device measurements. The fused evidence sources are dependent, and the reward is not ground-truth feedback.

\textit{Statistical validity.} Ten seeds in primary comparisons and five in secondary studies provide limited precision. Saturated scores, tied paired differences, and broad Holm families limit power. Bootstrap intervals are conditional on the observed seed sample, while sign-flip testing relies on the corresponding exchangeability/symmetry interpretation. Failure to reject is not equivalence. The component suite followed examination of the original outcomes and is reported as follow-up evidence, not a preregistered confirmatory study. Recovery is censored and should not be averaged as though every run recovered.

\textit{External validity.} Native Cooja motes, the unit-disk medium, CSMA, synthetic waypoint mobility, and one RPL implementation do not reproduce every low-power industrial network. Reference-paper parameter ambiguities prevent exact numerical replication claims. There is no external public intrusion dataset, physical deployment, independent research-group reproduction, authenticated adversary model, or target-MCU runtime study. Large-scale routing feasibility and subsequent attack comparisons remain unfinished. Intermittent forgers, contaminated calibration, TSCH/Orchestra, unseen attack families, federation, and lifelong retention require separate studies.

\section{Reproducibility and Research Integrity}\label{sec:repro}
The artifact contains firmware, the pinned Contiki patch and commit, scenario manifests and coordinates, mobility paths, selected binary definitions/hashes, compressed raw logs, parsers, statistical outputs, and figure generators. The original input manifest covers the 1960-run campaign. An independent replay of its log analysis matched all logged input hashes and eight numerical result files. The completed additional pipeline independently reparsed 740 compressed logs with zero core-result mismatches and exact agreement for four numerical tables. This establishes repeatable analysis of stored observations in the provided environment, not independent field replication.

From the project root, \texttt{make reproduce} regenerates the original analysis without Cooja. Running \texttt{analysis/reproduce\_journal\_extension.py} independently reparses the additional logs; the extension analyzer and plotter rebuild its tables and figures. The full-paper asset generator produces the new figures/tables and records their numerical input hashes. The supervisor guide distinguishes this fast path from hours of fresh simulation. CPU timing cannot be reproduced bit-for-bit on another host and must be reported with its environment. No public repository DOI or permanent archive accession is asserted; a reviewed public artifact release remains a submission task.

The experiments generate synthetic network traffic without private user data or live-network attacks. Failed scale controls and negative learning results are retained. Generative AI assisted manuscript drafting, code editing, and analysis organization; its outputs do not constitute experimental evidence. The author must review the technical statements, source attribution, implementation, and final disclosure against the selected journal policy. The bibliography retains the recorded DOI metadata audit; metadata matching does not imply that every full text was accessible or that every baseline implementation was independently certified.

\section{Conclusion and Future Work}\label{sec:conclusion}
This study examined concurrent DIS flooding and incorrect-parent DAO manipulation originating from the same compromised RPL nodes. AEF-RPL adds locally calibrated DAO-consistency evidence and selective handling to an independent DIS-FDM implementation, with a tabular Q-learning decision layer evaluated against a learning-disabled variant. Controlled simulations distinguish detection outcomes from the service retained by legitimate clients.

In the completed 30-node studies, AEF variants often preserve higher downward delivery than the DIS-focused comparators, while the route-repair ablation reveals a substantial descriptive service loss when that handling is removed. These benefits coexist with substantial mobile false isolation. Learning does not reduce false positives in the fresh-seed benign comparison and can produce repeated release/blacklist operations. Broadly adjusted extension tests also limit claims of statistical superiority. The defensible conclusion is a measured service--attribution trade-off, not uniform dominance or validated deployment readiness.

Large-network controls failed despite high conditional delivery among participating clients. Stack traces expose source-header fit rejections, neighbor-queue drops, and repeated DAO-triggered route repair; corrective validation remains separate from the reported defense results. Future work should complete those connectivity gates and independently seeded scale comparisons, measure constrained-device cost, and redesign adaptation using more reliable feedback before testing contaminated calibration, intermittent attacks, or additional MACs. The stored-log artifact supports inspection and regeneration of the present findings, including the negative results.
